\documentclass[12pt]{article}
\usepackage[utf8]{inputenc}
\usepackage[T1]{fontenc}
\usepackage{hyperref}
\title{Waigali Language Resources for K-12 Curriculum Development}
\author{Compiled by Claude Sonnet 5 (claude-sonnet-5) for Andr\'as Kornai}
\date{}
\begin{document}
\maketitle
\section{Introduction}
Waigali (ISO 639-3 wbk; also called Kalasha-ala) is spoken by roughly 10,000--12,000 people in the Waigal Valley, Nuristan, Afghanistan. Very little material exists. No primers, textbooks or school programmes were found in search results.
\section{Primary Educational Resources}
None found.
\section{Linguistic and Reference Materials}
\begin{itemize}
\item Morgenstierne, G. (1954), ``The Waigali Language'', \emph{Norsk Tidsskrift for Sprogvidenskap} 17: 146--219, cited in the Endangered Languages catalogue: \url{https://endangeredlanguages.com/node?page=389}
\item Glosbe community-built Waigali--English dictionary (informal, unverified entries): \url{https://glosbe.com/wbk/en}
\item English Wiktionary entry for Waigali (name and synonym Kalasha-ala): \url{https://en.wiktionary.com/wiki/Waigali}
\end{itemize}
\section{Religious and Community Resources}
None found.
\section{Digital and Audio Resources}
None found beyond the Glosbe pages above.
\section{Curriculum Development Context}
The language appears to have no school presence and no standardized orthography in the sources found. Any curriculum would have to begin from Morgenstierne's description and new community-led documentation.
\end{document}